\documentclass[12pt]{article}

\usepackage[a4paper, margin=2.5cm]{geometry}

\usepackage{circuitikz}

\usepackage{fontspec}
\usepackage{polyglossia}

\setmainlanguage{serbian}
\setmainfont{Times New Roman}

\title{Vežbe 1 - Ponavljanje osnovih koncepata elektrotehnike i uvod u LTSpice}
\date{}
\author{}
%\ctikzset{voltage/distance from node=0.25}
\tikzset{
    my voltage/.style={
        circuitikz/voltage/distance from node=0.75,
    }
}

\begin{document}

\maketitle

\section{Naponski razdelnik}

{Omogućava izračunavanje vrednosti pada napona na komponentama koje se nalaze u istoj grani, odnosno kroz koje protiče ista struja.}

\[
\begin{circuitikz}[american]
    \draw
        (0,0) node[ground] {}
        to[V, invert,l=$V_{in}$] (0,3)
        to[R, l=$R_1$, v_>=$V_{R_1}$, my voltage] (7,3)
        to[R, l=$R_2$, v=$V_{R_2}$, my voltage] (7,0)
        (7,0) node[ground] {};

    \draw[->, thick]
        (3,1) arc[start angle=180,end angle=0,radius=0.5]
        node[midway, below=0.2cm] {$K_1$};

    \draw[->, thick]
        (0,3) -- (2,3)
        node[midway, above] {$I$};
\end{circuitikz}
\]

{Na osnovu Omovog zakona dobijaju se izrazi}

\begin{equation}
    V_{R_1} = R_1 \cdot I
    \quad i \quad
    V_{R_2} = R_2 \cdot I.
    \label{eq:ohm}
\end{equation}

{Primenom drugog Kirhofovog zakona na konturu $K_1$ dobija se izraz za napon u kolu}

\begin{equation}
    -V_{in} + V_{R_1} + V_{R_2} = 0, odnosno
\end{equation}

\begin{equation}
    V_{in} = V_{R_1} + V_{R_2}.
    \label{eq:cont}
\end{equation}

{Nakon uvrstanja jednačine (\ref{eq:ohm}) u (\ref{eq:cont}) dobija se}

\begin{equation}
    V_{in} = R_1 \cdot I + R_2 \cdot I,
\end{equation}

{odakle se dobija izraz za struju u kolu}

\begin{equation}
    I = \frac{V_{in}}{R_1 + R_2}.
    \label{eq:current}
\end{equation}

{Gde se uvrštanjem jednačine (\ref{eq:current}) u izraz za napon na otpornicima (\ref{eq:ohm}) dobija}

\begin{equation}
    V_{R_1} = \frac{V_{in}}{R_1 + R_2} R_1 \quad i \quad V_{R_2} = \frac{V_{in}}{R_1 + R_2} R_2.
\end{equation}

{Naponski razdelnik deli ulazni napon na dva napona čije vrednosti zavise od otpronosti samih otpornika.}

{Struja u kolu ostaje ista dok su vrednosti napona različite za svaku komponentu.}

\section{Strujni razdelnik}

{Omogućava izračunavanje vrednosti struje kroz komponente koje imaju jednak pad napona.}

\[
\begin{circuitikz}[american]
    \draw
        (0,0) node[ground] {}
        to[I, l=$I_{S}$] (0,3)
        -- (3.5,3)
        to[R, l=$R_1$, v=$V_{R_1}$, my voltage] (3.5,0) node[ground] {}
        (3.5,3) -- (7, 3) to[R, l=$R_2$, v=$V_{R_2}$, my voltage]
        (7,0) to (7,0) node[ground] {};

    \draw[->, thick]
        (0,3) -- (2,3)
        node[midway, above] {$I_{in}$};
    \draw[->, thick]
        (3.5, 3) -- (3.5, 2.5)
        node[midway, right] {$I_1$};
    \draw[->, thick]
        (7,3) -- (7, 2.5)
        node[midway, right] {$I_2$};
    \draw
        (3.5,3) node[circ, label=above:$A$] {};

    \draw[->, thick]
        (1,0.75) arc[start angle=180,end angle=0,radius=0.5]
        node[midway, below=0.2cm] {$K_1$};

    \draw[->, thick]
        (5,0.75) arc[start angle=180,end angle=0,radius=0.5]
        node[midway, below=0.2cm] {$K_2$};
\end{circuitikz}
\]

{Treba primetiti da je pad napona na komponenatama u kolu jednak, zato što su one u paralelnoj vezi, odnosno važi da su $V_{R_1}$ i $V_{R_2}$ jednaki.}

\end{document}